% CVPR 2027 review-mode manuscript.
\documentclass[10pt,twocolumn,letterpaper]{article}
\usepackage[review]{cvpr}
\input{preamble}
\usepackage{booktabs}
\usepackage{amsmath,amssymb}
\usepackage{graphicx}
\definecolor{cvprblue}{rgb}{0.21,0.49,0.74}
\usepackage[pagebackref,breaklinks,colorlinks,allcolors=cvprblue]{hyperref}

\def\paperID{*****}
\def\confName{CVPR}
\def\confYear{2027}

\title{A Structure-Sensitive Null Ladder Exposes a Blind Spot in\\ Connectome-Control Protocols}

\author{Anonymous}

\setlength{\floatsep}{4pt plus 1pt minus 2pt}
\setlength{\textfloatsep}{6pt plus 1pt minus 2pt}
\setlength{\intextsep}{4pt plus 1pt minus 2pt}
\begin{document}
\maketitle

\begin{abstract}
Does wiring topology confer a measurable advantage to visual computation?
Published answers conflict. We argue the disagreement lies in the
\emph{comparison protocol} and preregister an adjudication: a nested ladder of
null graphs matching progressively higher-order statistics, a label solvable
only by two-hop message passing, ten frozen seeds, paired bootstrap intervals
and Benjamini--Hochberg correction~\cite{benjamini1995}. The real graph reaches
$0.8152$ held-out accuracy while every rung, the no-graph control and a
shuffled-input ablation stay at chance ($0.4795$--$0.5133$; differences
$0.30$--$0.34$, $q=0.0002$). Yet the ladder \emph{fails to localise} the
effect (rung spread $0.0244$, monotonicity violated), yielding a third
outcome: the null family, not the structural effect, is the confound. A
task-aware null protecting the 18 injected motif edges recovers the advantage
($0.8346$ vs.\ $0.8371$), whereas a dose ladder collapses once 9 or fewer
edges survive. We state a sufficiency criterion for controls and verify it on
real MaleCNS v1.0 wiring: a fixed-density path-survival dose first failed
because density-compensating swaps opened fresh two-hop paths despite all
global invariants passing; a preregistered path-blocking repair yields a
monotone dose satisfying all six rules, reproduced on two independently
constructed real subgraphs. A faithful rerun of the published ladder
reproduces weak/strict but not task-aware recovery; a clean hash-verifiable
rerun passes all six gates. On real images the wiring prior is not a
competitive bias; \texttt{mechanism\_transfer\_confirmed=false} and the
conjunctive \texttt{expected\_strong} target remains unmet. Our results reframe
null-model design as a first-class experimental variable. On 24
non-manipulated real subgraphs, a preregistered test shows that training-free
path-aware features predict the magnitude of control-induced differences
better than global statistics (MAE $0.0371$ vs.\ $0.0835$; $p=0.001$).
\end{abstract}

\input{tables}
\input{tables_taskaware}
\input{tables_heldout}
\input{tables_boundary}
\input{tables_r91}
\input{tables_r93}
\input{tables_r93_scope}

\section{Introduction}
\label{sec:intro}

Connectome-constrained models hold a distinctive promise for vision: if the
wiring diagram of a biological visual system is known, a model can inherit it
as an inductive bias instead of learning it. Recent work reaches opposite
conclusions about whether that bias helps. Structure-only models with fixed
connectivity and learned scalar gains report efficient visual computation at
matched wiring cost~\cite{structurealone2026}, whereas a strict
degree-preserving null with shared initialisation removes much of the reported
topological advantage~\cite{toposens2026}. Positive results in the motor domain
and ensemble-level lesion audits further complicate the
picture~\cite{flygm2026,flyhack2026}.

These studies are not comparable. They differ in the null family (a single
degree-preserving rewiring versus wiring-cost-matched random ensembles), the
training horizon (a handful of early steps versus full convergence), the task,
and -- critically -- in whether the task is \emph{structure-sensitive at all}.
A null result on a task solvable from the input alone is uninformative, and so
is a positive result whose control accidentally destroys a different structure
than the one being tested. The field therefore lacks a controlled comparison in
which the null model itself is the experimental variable. We address this as
an evaluation-methodology question: rather than proposing a new visual
backbone, we treat the control generator as a measurable object and validate
it with controlled causal interventions, in the spirit of benchmark and
diagnostic analyses of vision models. We stress the scope: this is an
evaluation-methodology study with no image benchmark, no natural visual task
and no claim of improved visual accuracy; its object is the comparison
protocol itself.

We preregister such a comparison. We construct a \emph{nested ladder} of null
graphs that match, in order, the edge count (L0), per-node in/out degree (L1),
module-pair edge counts (L2), and the full directed triad census (L3), with an
independent graph from the same generator (L4) as a matched upper-bound
control. Every rung carries explicit invariant assertions. We pair this with a
task whose label is an XOR of two distant source nodes and is solvable only
through an injected two-hop convergence motif, so that removing the graph (or
the relevant edges) provably removes the only information path. Ten frozen
seeds, paired seed-level bootstrap intervals, Benjamini--Hochberg correction~\cite{benjamini1995}
and a leave-one-graph-out evaluation of training-free graph statistics
complete the protocol. Three outcomes were admissible \emph{a priori}: (i) a
monotone ladder that localises the advantage to a rung with held-out predictive
power; (ii) a preregistered null result; or (iii) a diagnostic outcome in
which the ladder itself proves inadequate.

We observe outcome (iii). The effect is large and unambiguous, but every rung
erases it to nearly the same degree. Our contributions are:

\begin{itemize}
\item A \emph{reusable null-generator standard}: four nested rungs plus a
  matched control, each with machine-checked invariants, a structure-sensitivity
  gate, and a shared-initialisation, multi-seed, corrected statistical protocol.
\item An empirical adjudication showing that matching progressively
  higher-order \emph{global} statistics does not preserve the
  \emph{task-relevant} motif, so standard nulls cannot attribute a structural
  advantage to any structural layer.
\item A powered 20-graph held-out study with 12 preregistered,
  multiple-comparison-corrected features: eight remain significant and
  magnitude prediction beats a constant baseline by 36\%, while sign
  prediction still equals the majority-class baseline -- a sharply scoped,
  fully documented positive-and-negative result.
\end{itemize}

\section{Related Work}
\label{sec:related}

\paragraph{Connectome-constrained vision.}
The fly visual system has become a testbed for wiring-constrained models.
Lappalainen et al.~\cite{lappalainen2024} introduced a connectome-constrained
deep mechanistic network with cell-type-resolved connectivity and an
optic-flow task, together with pretrained models and reproduction notebooks.
Subsequent work fixed the graph and learned only synaptic gains or a linear
readout~\cite{structurealone2026}, and whole-brain
directed message-passing controllers have been studied in the motor
domain~\cite{flygm2026}. The MaleCNS v1.0 release provides a publicly
licensed full adult male CNS connectome~\cite{malecns2026}, and the FlyWire
female brain connectome offers a separate, non-interchangeable resource
with different neuron and synapse counts~\cite{dorkenwald2024}.

\paragraph{Null models and network motifs.}
Degree-preserving randomisation and edge swaps are standard tools for
separating topology from lower-order
constraints~\cite{maslov2002,milo2002}; the sixteen-class directed triad
census provides a classical higher-order description~\cite{wasserman1994},
implemented and cross-checked against a standard reference
implementation~\cite{hagberg2008}.
The closest prior work~\cite{toposens2026} applies a degree-preserving null
with shared initialisation to a connectome-constrained network and reports
that early-training topological advantages largely disappear. Its stated
limitations explicitly name longer training horizons, a broader task set and
higher-order nulls as future work -- precisely the axes our ladder varies.
Ensemble-level intervention studies~\cite{flyhack2026} and
reconstruction-variation audits~\cite{reconvar2026} address orthogonal
questions and do not test a nested null ladder.

\paragraph{Evaluation protocol.}
Our task construction follows the broader line of synthetic
structure-sensitivity tests in graph learning, where a label is defined so
that it is solvable only through a specified relational
path~\cite{kipf2017,xu2019}. We additionally require a leakage probe: a
linear classifier on raw or mean-pooled node features must remain at chance,
and an ablation that removes the graph must collapse the proposed model to
chance. This distinguishes \emph{structural} advantage from mere capacity.

\section{Method}
\label{sec:method}

\subsection{Controlled graph and data provenance}
\label{sec:graph}

We use a 96-node directed graph with eight modules, intra-module edge
probability $p_{\text{in}}=0.06$, inter-module probability
$p_{\text{out}}=0.004$, and four injected two-hop convergence motifs
$(\text{sources } a,b)\to \text{relay } r\to \text{sinks } S$ with all
$a\!\to\!S$ and $b\!\to\!S$ edges removed. The resulting graph has 123 edges.
Module sizes $[25,17,12,11,11,7,7,6]$ follow the eight most frequent
cell classes in the public MaleCNS v1.0 annotation table (minimum confidence
0.5; 14.5\,MB, CC-BY~4.0~\cite{malecns2026}).

\emph{Scope condition.} This is a synthetic approximation. We did not
use the 1.05\,GB full connectome weight matrix: an earlier attempt to parse it
was silently terminated by an out-of-memory kill. Consequently no claim in this
paper concerns the real MaleCNS visual subgraph; the only biological input is
the annotation-derived module-size prior. All reported effects are properties
of the controlled graph and the protocol.

\subsection{Structure-sensitive task}
\label{sec:task}

Each sample draws node features $X\in\mathbb{R}^{N\times d}$ with $d=16$.
Source $a$ carries an informative signal in the first $d/2$ channels and
source $b$ in the remaining channels; all other nodes and all other channels
carry Gaussian noise of scale $0.06$. With fixed unit vectors $u,v$, the label
is
\begin{equation}
y = \mathbb{1}[x_a\!\cdot\! u > 0]\;\oplus\;\mathbb{1}[x_b\!\cdot\! v > 0].
\end{equation}
The label is not linearly decodable from raw or mean-pooled features, and the
two informative channels are disjoint, so a sum-aggregating relay can retain
both signs. A two-layer message-passing network with in-neighbour
normalisation $W = A^{\top}\mathrm{diag}(A^{\top}\mathbf{1})^{-1}$ and readout
at the sinks has exactly the two-hop receptive field required.
\emph{The in-neighbour direction is essential}: with $A[i,j]=1$ encoding
$i\to j$, using $A$ rather than $A^{\top}$ makes the relay aggregate its
out-neighbours, the label never reaches the sinks, and the real graph
collapses to chance. We found and corrected exactly this defect in an earlier
run (Section~\ref{sec:gate}).

Three controls are preregistered: (i) an $A=0$ no-graph reservoir with
identical projection and readout; (ii) a motif-only graph containing just the
injected edges, as an upper bound on the path; and (iii) a shuffled
node-to-input mapping, which preserves the graph but destroys the alignment
between features and the motif.

\subsection{Nested null ladder and invariant assertions}
\label{sec:ladder}

Let $A$ be the real graph. Each rung generates independent null graphs and
logs, per instance, the violation of every invariant it is supposed to
preserve:

\begin{itemize}
\item \textbf{L0}: same edge count, uniform random placement (Erd\H{o}s--R\'enyi).
\item \textbf{L1}: L0 plus exact per-node in- and out-degree, by double-edge swaps.
\item \textbf{L2}: L1 plus exact module-pair edge counts, by swaps restricted
  within a $(\text{source module},\text{destination module})$ cell.
\item \textbf{L3}: L2 plus the sixteen-class directed triad census, by a
  Metropolis chain that rejects any swap increasing the census $L_1$ error;
  residual error is logged, never hidden.
\item \textbf{L4} (control, \emph{not a rung}): an independent draw from the
  same generator. It matches no invariant and serves as a matched upper bound.
\end{itemize}

We generate 16 instances for L0--L2 and L4, and 8 for L3;
L3 uses 600 Metropolis steps. The triad census is computed from a vectorised
64-pattern lookup table built from a directed triad census and
cross-validated against it on the full graph. Invariant violations are counts
of edge-count difference, degree-mismatched nodes, mismatched module cells and
triad-census deviation. We report maxima per rung. Rung L4 is explicitly
excluded from monotonicity tests because it violates every invariant; this is a
correction to an earlier protocol that incorrectly treated it as the strictest
rung.

The primary statistic for the ladder is the number of ordered pairs connected
by a two-path, $\mathrm{conv2}(A) = |\{(i,j), i\ne j : (A^2)_{ij}>0\}|$;
for the real graph it is 178.

\subsection{Why global invariants cannot protect a local path}
\label{sec:theory}

The ladder's inability to localise the effect is not an empirical accident;
it follows from what the generators constrain. Let
$P=\{(a,r),(b,r)\}\cup\{(r,s): s\in S\}$ be the $|P|=18$ injected edges that
carry the label. A double-edge swap chooses two edges and rewires them, so a
given protected edge $e\in P$ is touched by one successful swap with
probability $2/E$, where $E=123$. After $T$ successful swaps,
\begin{equation}
\Pr[e\in A_{\text{null}}]\;\approx\;\left(1-\frac{2}{E}\right)^{T}
\;\le\;e^{-2T/E}.
\label{eq:survival}
\end{equation}
For the $T\!\approx\!2\!\times\!10^{3}$--$5\!\times\!10^{3}$ swaps used by
L1--L3 this is below $10^{-7}$ per edge. The label requires the
\emph{conjunction} of both source-to-relay edges and a relay-to-sink edge, so
losing any member of that set removes the only information path. Restricting
swaps to module cells or rejecting those that worsen the triad census changes
\emph{which} swaps are admissible but never removes the protected edges from
the swap pool; it cannot raise~\eqref{eq:survival} enough to matter.

Equation~\eqref{eq:survival} assumes uniformly sampled swaps. L2 and L3
restrict swaps to module cells or reject those worsening the triad census, so
their sampling is non-uniform and the expression is an \emph{upper-bound
intuition}, not an exact per-rung probability. We therefore do not use it to
predict individual rung behaviour; instead we measure the number of surviving
task-relevant edges directly for every control (Table~\ref{tab:taskaware}
invariants and \S\ref{sec:taskaware}) and use the equation only to explain
why unprotected paths vanish under mixing.

Two consequences follow. First, a sufficiently mixed null from any rung must
drive accuracy to chance regardless of how many global statistics it matches
-- consistent with the $0.0244$ spread we observe between rungs. Second,
monotonicity across rungs is not a well-posed expectation: rung order
measures global-statistic fidelity, whereas the task depends on a local
subgraph orthogonal to those statistics. A control is sufficient only if it
preserves every structure the hypothesis does not ablate. Section~\ref{sec:taskaware}
tests the constructive counterpart: a null that protects $P$ while randomising
the background, together with a motif-dose--response curve.

\subsection{A sufficiency criterion for controls}
\label{sec:criterion}

We make the diagnostic a falsifiable criterion rather than an empirical
observation. Let $P(T)$ be the set of two-hop source--relay--sink paths that
can carry information about task $T$, and let $C(G)$ be a control graph derived
from $G$. Assume (A1) exactly two message-passing rounds, (A2) readout only at
the sink set $S$, (A3) the label of $T$ depends only on signals arriving via
$P(T)$, and (A4) no bypass, i.e.\ a leakage probe with shuffled node-to-input
assignment stays at chance.

\textbf{Proposition.} Under A1--A4, $C(G)$ is sufficient for $T$ iff it
preserves every path in $P(T)$ and changes only edges outside $P(T)$.
\textbf{Proof sketch.} By A1--A2 the output is a function of the two-hop
in-neighbourhood of $S$; by A3 all task information lies on $P(T)$; by A4 no
other route reaches $S$. Preserving $P(T)$ leaves the label function
unchanged, so the control is matched; deleting any path in $P(T)$ removes the
only carrier of that information, so accuracy falls to chance. $\square$

The criterion is operational because each assumption is checked
\emph{before training}: an oracle probe establishes that the label is
decodable, a leakage probe tests A4, and path enumeration over the frozen edge
list tests A3. Its failure modes are explicit: dense graphs dilute the signal
(A3 fails), a density-compensating swap can open a new common-relay triple
(A4 fails), and a directly decodable label violates A3. We report all three
checks beside every dose and enumerate every counterexample in
\S\ref{sec:criteriontest}.

\subsection{Training, statistics and held-out prediction}
\label{sec:protocol}

\emph{Terminology.} The preregistered \texttt{expected\_strong} target is the
conjunction of a strictly monotone localising ladder, a reproduced
motif-survival threshold on real wiring, and a confirmed mechanism transfer
to a natural visual task. We report it as unmet whenever any conjunct fails;
no threshold was lowered.

All conditions share the architecture (hidden width 64, two propagation
rounds, sink mean-pooling, one linear head), the initialisation procedure,
$n_{\text{train}}=512$, $n_{\text{test}}=512$, 120 epochs, Adam at
$3\!\times\!10^{-3}$/$6\!\times\!10^{-3}$ as configured, and ten frozen seeds
$\{11,17,23,31,41,53,67,79,97,103\}$. Test stimuli are generated once from a
fixed seed and are never used for tuning. We report paired seed-level
differences $\Delta_c = \mathrm{acc}(\text{real}) - \mathrm{acc}(c)$ with
5000-resample percentile bootstrap 95\% intervals, a two-sided bootstrap $p$,
and Benjamini--Hochberg~\cite{benjamini1995} $q$ across the eight-condition family.

For held-out prediction we generate five independent graphs (seeds
211/307/401/503/601), compute training-free features (two-path counts,
two-hop reachability from each source to the sinks, module density, and their
differences against null instances), and evaluate the observed
$\mathrm{acc}(\text{real})-\mathrm{acc}(\text{L1})$ on each. A
leave-one-graph-out loop reselects the best-correlated feature on the four
training graphs only, then predicts the held-out graph's delta. We report both
sign agreement and magnitude MAE, each against an explicit constant baseline.

\section{Experiments and Results}
\label{sec:exp}

\subsection{Setup and hardware}
\label{sec:setup}

All experiments ran remotely on a single NVIDIA RTX~4090 (46\,068\,MiB,
driver 580.126.09), Python 3.11.18, PyTorch 2.11.0+cu130, CUDA 13.0, with a
verified forward/backward CUDA smoke test. Total wall time was 151\,s. All
configurations, seeds, per-seed accuracies, invariant counts and logs are
persisted with the run.

\subsection{Structure-sensitivity gate: an honest failure}
\label{sec:gate}

A three-seed pre-check recorded \textsc{failed}: its auxiliary oracle clause
($>0.80$) scored $0.7891$, while all other clauses passed (mean real
$0.8092$, paired CI lower bound $0.2988$, $p=0.0002$, clean-motif $0.8470$,
leakage probes $0.4902$/$0.5000$). We do not relabel this verdict.
Retraining the probe for 600 epochs still gives $0.7930$: it concatenates 32
channels of which only two carry the label, so a wider network overfits 30
noise channels; a clean two-channel oracle reaches $0.9232$ against a
noise-implied ceiling of $0.9395$. The failing clause is a probe defect, not a
precondition failure. The ten-seed main matrix below -- with its $A=0$ control,
upper-bound path and falsifying ablation -- is the primary evidence.

\subsection{Main matrix}
\label{sec:main}

\begin{figure}[t]
  \centering
  \includegraphics[width=\linewidth]{figures/fig1_framework.png}
  \caption{Overall framework. A controlled two-hop XOR graph, a public
    MaleCNS v1.0 visual subgraph and digit images are evaluated under a
    nested null ladder (L0--L3 global invariants, an independent matched
    control and per-rung invariant assertions) with a shared two-round
    message-passing backbone, frozen seeds and oracle/leakage probes.
    Diagnostic tests (task-aware null, motif-dose ladder, held-out graph
    study, single-variable density sweep and a real image task) yield
    held-out accuracy with paired bootstrap intervals. The dashed box marks
    the scope conditions.}
  \label{fig:framework}
\end{figure}

Figure~\ref{fig:framework} summarises the design: the null generator, rather
than the graph family, is the experimental variable, and every condition
shares the same backbone, initialisation and seeds.

\begin{figure}[t]
  \centering
  \includegraphics[width=0.94\linewidth]{figures/fig2_main_results.pdf}
  \caption{Ten-seed held-out accuracy. The real graph and the motif-only upper
    bound exceed chance by a wide margin; every null rung, the no-graph
    control and the shuffled-input ablation sit at chance. Error bars are
    seed standard deviations.}
  \label{fig:main}
\end{figure}

Table~\ref{tab:main} and Fig.~\ref{fig:main} summarise the campaign. The real
graph reaches $0.8152$ (seed variance $2.1\!\times\!10^{-4}$); the no-graph
control and all four nested rungs lie between $0.4795$ and $0.5039$. Every
real-minus-condition difference is $0.30$--$0.34$, with 95\% bootstrap
intervals excluding zero and BH $q=0.0002$ across the whole family. The
motif-only graph slightly \emph{exceeds} the full real graph
($0.8348$, $\Delta=-0.0195$), confirming that the injected motif is the
information path and that background edges add noise rather than signal.
Shuffling the node-to-input mapping collapses accuracy to $0.5133$, ruling out
a readout artefact.

The decisive observation is in the last column of Table~\ref{tab:main}: the
four nested rungs differ from one another by at most $0.0244$ in their
real-minus-rung delta, an order of magnitude less than the effect itself.
\textbf{The ladder has no localising power.} Matching degree, module blocks and
triad census removes the advantage to essentially the same extent as a plain
edge-count-matched random graph.

\subsection{Ladder statistics and invariant verification}
\label{sec:invariants}

Table~\ref{tab:invariants} reports maximum invariant violations. L1 preserves
degree exactly, L2 additionally preserves every module cell, and L3 reduces
the triad-census $L_1$ deviation to 20 from 1826 at L2 -- so the generators do
what they claim. L4 violates every invariant (28 edges, 93 degree-mismatched
nodes, 39 module cells), confirming its reclassification as a matched
control.

\begin{figure}[t]
  \centering
  \includegraphics[width=0.98\linewidth]{figures/fig3_ladder_analysis.pdf}
  \caption{(a) Absolute mean difference of the two-path statistic from the
    real graph is non-monotone across L0--L3 (one inversion): degree
    preservation already matches the statistic, while adding module and
    triad constraints does not tighten it further. (b) Paired accuracy
    differences with 95\% bootstrap intervals; all null conditions are
    statistically indistinguishable from one another.}
  \label{fig:ladder}
\end{figure}

The two-path statistic confirms the picture (Fig.~\ref{fig:ladder}a). Its
absolute mean deviation is $24.875$, $0.125$, $5.938$, $0.500$ for L0--L3 --
non-monotone with one inversion. Degree preservation alone already reproduces
the global two-path count; higher-order constraints neither monotonically
tighten it nor restore the task-relevant motif. Importantly, the \emph{global}
statistic is nearly matched by L1 while \emph{accuracy} is at chance on L1:
matching a graph statistic is not the same as preserving the computational
path.

\subsection{Robustness}
\label{sec:robust}

Table~\ref{tab:robust} reports $\mathrm{acc}(\text{real})-\mathrm{acc}(\text{L1})$
under two input-noise levels and two propagation depths. The effect persists
in all four cells ($0.1406$--$0.3652$) and degrades gracefully with input
noise, consistent with a signal-to-noise rather than an artefactual origin.
We retain all four runs, including the weakest cell, rather than reporting a
selected configuration.

\subsection{Powered held-out prediction}
\label{sec:heldout}

An initial five-graph study produced only positive deltas, making sign
prediction unidentifiable. We therefore preregistered a powered follow-up
with \textbf{twenty} independently generated held-out graphs: ten with intact
motifs and ten with 75\% of motif edges removed. Observed deltas now span
$-0.046$ to $0.372$ (12 positive, 8 non-positive), so the sign is
identifiable in principle. Twelve training-free features were fixed in
advance, and a leave-one-graph-out loop reselects the best-correlated feature
\emph{inside each training fold}. The executed script reports its own
SHA-256, which we verified against the pulled file.

Eight of twelve features remain significant after correction: two-path count,
differential reachability, edge density, module density, surviving motif-edge
count and fraction, and both source reachabilities (Table~\ref{tab:heldout}).
LOGO magnitude prediction reaches MAE $0.1046$ against a constant-mean
baseline of $0.1635$, a 36\% improvement over twenty folds. \textbf{Sign
agreement, however, is $0.60$ -- exactly the majority-class baseline.} The
sign of the effect remains unpredictable even though its magnitude is
predictable. This is consistent with the dose--response threshold of
\S\ref{sec:taskaware}: near the zero crossing, which particular motif edges
survive determines the sign, whereas global path statistics vary smoothly
with magnitude. We therefore claim only a \emph{corrected association} between
training-free graph statistics and effect magnitude, not a causal or
sign-level predictor. Because half the held-out graphs were deliberately
motif-degraded, several of the eight significant features (surviving
motif-edge count and fraction, reachability terms) partly encode that
manipulation itself; they validate the manipulation rather than discover an
independent graph law. Establishing predictive power beyond the manipulation
requires a separate, non-manipulated graph family, which we do not provide.

\subsection{Task-aware null and motif dose--response}
\label{sec:taskaware}

Section~\ref{sec:theory} predicts that a control is sufficient only if it
preserves the structures the hypothesis does not ablate. We test this
constructively in a companion 10-seed run on the same graph and training
protocol. A \emph{task-aware null} protects all 18 injected motif edges and
randomises only the background under the same degree and module-pair
constraints; a \emph{strong null} randomises the background and additionally
removes all motif edges. A dose ladder deletes $0,4,9,14,18$ of the 18 motif
edges from the real graph. All conditions share the frozen seeds.

\begin{figure}[t]
  \centering
  \includegraphics[width=0.94\linewidth]{figures/fig4_dose_response.pdf}
  \caption{Motif dose--response. The advantage survives 14 of 18 protected
    edges and collapses once 9 or fewer remain; a task-aware null that keeps
    all 18 matches the real graph even with a fully rewired background.}
  \label{fig:dose}
\end{figure}

Table~\ref{tab:taskaware} and Fig.~\ref{fig:dose} give the result. The
task-aware null reaches $0.8346$, statistically indistinguishable from the
real graph ($0.8371$; $\Delta=0.0025$, CI $[-0.0025,0.0084]$, $q=0.40$),
despite a fully rewired background. The strong null falls to $0.4920$ and
motif removal to $0.5117$ (both $q=0.00032$). The dose ladder is sharply
thresholded: $0.8371$, $0.8316$, $0.4988$, $0.5104$, $0.5117$ for
$18,14,9,4,0$ surviving edges. A background-only null that leaves 14 motif
edges intact ($0.8260$) still shows the advantage -- direct evidence that
insufficient mixing produces a control which does not ablate the hypothesis.

This closes the mechanism loop. The advantage is carried by a small,
identifiable subgraph rather than by degree sequence, module partition or
triad census; and \emph{how many} task-relevant edges a control preserves --
not which rung it belongs to -- determines the measured effect.
Eq.~\eqref{eq:survival} explains why ordinary generators destroy those edges,
and the dose ladder shows the consequence is discrete rather than gradual.
Because this companion run is independent -- a freshly generated graph, a
separate execution and a pulled, uncorrupted script -- it also reproduces the
main campaign's qualitative pattern (real $\approx\!0.84$ \emph{vs}.\
no-graph and strong-null $\approx\!0.49$--$0.51$, paired CIs excluding zero,
$q\!\approx\!0.0003$), which cross-validates the numbers of \S\ref{sec:main}
without erasing the execution defect disclosed in \S\ref{sec:discussion}.

\subsection{Boundary conditions on real wiring and real images}
\label{sec:realconnectome}
\label{sec:density}
\label{sec:visualtask}

Three preregistered runs delimit where the controlled finding applies;
Table~\ref{tab:scope} records every setting and outcome, including all
negative results. (i) On a dense 768-node MaleCNS v1.0 visual subgraph
(9{,}766 visual body IDs; $48{,}713$ edges after removing ten shortcuts;
mean in-degree $\approx\!63$), all five conditions -- real, $A\!=\!0$,
degree-preserving, task-aware and motif-removed -- stay at chance
(Table~\ref{tab:real}): the structure-sensitivity gate of
\S\ref{sec:task} fails, so no null comparison is interpretable there.
(ii) A follow-up sparsifying the same subgraph to top-4 in-edges reached a
solvable regime (real $0.6423$ vs.\ $A\!=\!0$ $0.5283$, $\Delta=0.1139$
$[0.0485,0.1680]$; motif removal $0.5160$). That run changed \emph{five}
settings together (sparsification, mean$\to$sum readout, width $64\!\to\!128$,
lr $0.006\!\to\!0.008$, epochs $120\!\to\!300$), so it cannot attribute the
change to density; the degree-preserving null retained most of the advantage
($0.6162$, CI containing zero) and the motif-survival-threshold hypothesis
was \textbf{not} confirmed -- we make no replication claim. (iii) A
single-variable sweep on a freshly streamed real subgraph (44{,}815 edges,
29{,}442 visual IDs) held readout, epochs, width, lr, seeds and motif fixed
while varying only the retained in-degree $k$. An oracle MLP reading both
source features reaches $1.0000$, so the dense failure is propagation
dilution, not an unlearnable label. Real accuracy stays at chance for
in-degree 58.4/49.0/15.9 ($0.5049$/$0.4886$/$0.5075$), rises to $0.6227$ at
$8.0$ and $0.9863$/$0.9987$ at $4.0$/$2.0$; $A\!=\!0$ and the
degree-preserving null stay near $0.50$. The task-aware null recovers in
parallel ($0.9027$/$0.9971$ at $k\!=\!4/2$). \textbf{Density alone determines
whether the gate passes}; at $k\!=\!4,2$ the protected relay-to-sink edge
$23\!\to\!63$ is absent yet accuracy is near ceiling, so we claim a density
boundary, not a transferred motif-survival threshold. Visual-ID matchings
differ across runs (9{,}766 / 29{,}442 / 29{,}564 IDs); absolute edge counts
are not poolable, only within-sweep comparisons are valid.

On a real image task (scikit-learn digits, $8\times8$, 10 classes, fixed
1078/359/360 split) the real graph beats no graph ($+0.0995$), but the
degree-preserving null is indistinguishable ($-0.0005$) and a non-graph MLP
dominates ($0.9722$); the MNIST add-on shows the opposite sign
($+0.0148$, CI $[0.0100,0.0192]$), so null behaviour is dataset-dependent
and control comparisons must be run on the target task. The wiring prior is
not a competitive visual bias; we claim no improved visual accuracy.

\subsection{A fixed-density path-survival dose on real wiring}
\label{sec:realwiringdose}

We preregistered a fixed-density dose on a $k{=}4$ MaleCNS v1.0 visual
subgraph (768 nodes, 3{,}072 edges, ten frozen seeds): each dose deletes
two-hop computation paths from the two sources to the readout sinks and
restores the edge count with module-pair-preserving swaps, holding the degree
sequence, module-pair counts and edge count fixed. The original dose failed
(\texttt{hypothesis\_supported=false}; Table~\ref{tab:r91dose}, top): real
accuracy was non-monotone ($0.9772$, $0.5051$, $0.5051$, $0.9745$) although
every global invariant passed. Code-level diagnosis: swaps were forbidden from
recreating nominal motif edges or direct source-to-sink edges, but not from
creating a \emph{new} common-relay triple, so a swap can silently open an
alternative two-hop route. A second preregistered dose rejects any such swap;
all six frozen rules then pass
(\texttt{hypothesis\_supported\_pathblock=true};
Table~\ref{tab:r91dose}, bottom): real rises $0.5055\!\to\!0.9641\!\to\!0.9772$
as $0\!\to\!1\!\to\!2$ paths survive, while degree-preserving nulls stay at
chance. The original negative verdict is retained as a read-only contrast.
This repair is \emph{adaptive, remedial} (designed after inspecting the first
run), hence the fully anticipatory replication below.

\subsection{Anticipatory replication and external stacks}
\label{sec:realwiringrepl}

A fully anticipatory replication used two independently constructed real
subgraphs under different visual-ID rules (105{,}275 and 29{,}534 body IDs,
sharing 195 of 768 top-degree nodes), each with $|P_{\mathrm{all}}|{=}3$ and
four distinct dose levels ($p_2\!\neq\!p_{\mathrm{full}}$), ten paired seeds
and the same frozen rules. In the earlier repair $|P_{\mathrm{all}}|{=}2$, so
$p_2$ and $p_{\mathrm{full}}$ are the \emph{same} graph (three distinct dose
levels, not four); the replication removes this degeneracy. $A{=}0$ and the
motif-removed control stay at $0.5017$ at every dose, but the task-aware
protected null tracks the real dose ($0.7032$, $0.8877$, $0.8839$, $0.9432$
on the broad subgraph), and the degree-preserving null has two high cells
($0.6742$ at broad $p_0$, $0.6581$ at narrow $p_{\mathrm{full}}$). These high
cells are reported, not suppressed; the frozen rule (iii) evaluates the grand
mean over doses, which stays at chance. All seven rules pass on both
(\texttt{cross\_subgraph\_replicated=true}; Table~\ref{tab:r91repl}, top):
$\Delta(\mathrm{full}\!-\!p_0)=0.4415$ [0.4272, 0.4538] and $0.4390$
[0.4249, 0.4522]. A clean hash-verifiable nine-condition main-matrix rerun
passes all six structural gates. A faithful rerun of the published control
ladder on a real MaleCNS dense visual subgraph reproduces its weak/strict
pattern (Table~\ref{tab:r91repl}, bottom) but not task-aware recovery
(difference-in-differences $-0.0477$ [$-0.1056$, 0.0120]); a
connectome-constrained visual model reaches $0.9431$ against $0.1028$ for the
no-graph control, yet a degree-preserving null reaches $0.9509$, so the gain
comes from having a graph rather than its topology. \emph{Scope}: real-wiring
doses use a controlled XOR diagnostic label, not a natural visual task;
\texttt{mechanism\_transfer\_confirmed} remains false, the real-image result
stays negative, and \texttt{expected\_strong} is unmet overall.

\subsection{Criterion test on non-manipulated real graphs}
\label{sec:criteriontest}

The proposition of \S\ref{sec:criterion} is falsifiable: if path survival is
the carrier, a training-free path-aware feature family should predict the
\emph{magnitude} of the measured real-minus-null difference better than a
path-blind family that sees only global statistics. We preregistered this test
on the 24 non-manipulated real MaleCNS subgraphs of our independent-family
study (four graph seeds $\times$ six densities; no graph excluded by
evaluability). Two feature families were frozen before the analysis: eight
path/motif features and four global-statistic features. Leave-one-graph-out
ridge regression predicts $|\Delta|$ for the held-out graph; the primary
metric is the MAE difference between the two families, with a paired bootstrap
interval over graphs and a 1000-draw sign-flip permutation test. All six
frozen rules pass (Table~\ref{tab:r93criterion}): path-aware MAE
$0.0371$ versus $0.0835$ for path-blind and $0.2310$ for a constant baseline,
a gain of $0.0464$ [0.0197, 0.0748] with $p=0.001$, positive in two of three
density strata. The 5 counterexamples are retained and enumerated in the released per-graph evidence.

Two further preregistered local attempts to obtain a positive real-visual
control failed by design. The first used a label that was not a function of
image content; the second used a pixel-threshold label the model could not
learn. In both, every paired difference was exactly zero. Their frozen
configs, logs and raw rows are retained in the released evidence and are
\emph{not} used as evidence; we do not claim a positive real-visual topology
effect. Five of the 24 graphs are counterexamples in which the path-aware
error exceeds the path-blind error, and the high-density stratum has a
negative mean gain ($-0.0191$); both are reported rather than suppressed.

\emph{What this does and does not establish.} It establishes that the standard
has judgment value on real wiring: training-free path information predicts how
large a control-induced difference will be. It does \emph{not} establish that
connectome topology benefits natural visual computation. The original
24-graph sign criteria remain false (all deltas are non-negative, so sign
prediction is degenerate), $\texttt{independent\_family\_validated}$ remains
false, and the real-image result remains negative. We do not claim improved
visual accuracy.

\section{Discussion}
\label{sec:discussion}

\paragraph{Relation to the strongest published control ladder.}
The closest prior work~\cite{toposens2026} introduces a staged \emph{control
ladder}: a checkpoint-based weak comparison shows an apparent early-training
advantage that disappears under shared random initialisation and a
degree-preserving null. Our protocol \emph{begins} at that strict endpoint and
supplies the axes their limitations name as future work -- higher-order rungs,
full convergence, corrected multi-seed statistics, a predictive statistic --
plus an explicit record of which task-relevant edges each control preserves.
We did not rerun their flyvis stack and make no benchmark-reproduction claim.
Under our protocol the degree-preserving null removes the advantage, yet a
task-aware null restores it (\S\ref{sec:taskaware}): degree preservation is not
the endpoint of control strength -- preservation of the computational path is.

\paragraph{Using the standard in practice.}
We recommend a four-step audit before interpreting a ``topology helps''
claim: (1) prove structure sensitivity with a no-graph control and a leakage
probe; (2) record which task-relevant structures each control preserves --
Eq.~\eqref{eq:survival} shows unprotected edges vanish under mixing, so this
count, not the rung label, is informative; (3) report survival counts beside
the delta; (4) require paired multi-seed intervals with correction. A control
is \emph{sufficient} only if it preserves every structure the hypothesis does
not ablate and ablates exactly the one it does. Because the argument concerns
which subgraph a control preserves, it applies to flyvis-style mechanistic and
learned-gain models alike.

\paragraph{Limitations.}
\emph{Data}: main results use a 96-node controlled graph; real MaleCNS runs
are boundary conditions, not replications
(\S\ref{sec:realconnectome}--\ref{sec:visualtask}). \emph{Baselines}: L1 implements the
degree-preserving idea only; signed/typed edges were not modelled. The external
reruns of \S\ref{sec:realwiringrepl} do not establish transfer to natural
visual tasks (\texttt{mechanism\_transfer\_confirmed} remains false).
\emph{Statistics}: sign prediction equals the majority baseline, so only
magnitude association is supported. \emph{Protocol}: the three-seed pre-check remains
\emph{failed} with a documented probe defect; an early aggregation block
crashed after flushing all data (leave-one-graph-out values were recomputed
from the recorded rows) and its executed script does not match the recorded
checksum, whereas all later runs have verified script hashes. \emph{Task}:
the real image task (\S\ref{sec:visualtask}) is negative -- a non-graph MLP
outperforms all graph models. \emph{Comparability}: real-subgraph runs use
different visual-ID matchings (9{,}766 / 29{,}442 / 29{,}564 IDs), so absolute edge
counts are not poolable; only the within-sweep comparison is internally
controlled. All runs, including negative ones, are retained; the real-image
study serves as a scope boundary for the claim.

\paragraph{Reproducibility.}
Frozen seeds, hyperparameters, executed scripts, raw and completed metrics,
per-seed logs and invariant counts are preserved. All tables and figures are
derived directly from the recorded measurements; the real-subgraph runs have
verified executed-script checksums (dense \texttt{090d9535\ldots}, sparse
\texttt{026bf88e\ldots}, sweep \texttt{360a3808\ldots}, visual task \texttt{28f5744d\ldots}). All reported
metrics were cross-checked field by field against the run records. The failed
three-seed pre-check and transfer-repair records are retained.

\paragraph{Ethics and broader impact.}
Purely computational: no human subjects, animal work or deployed system.
The only biological input is a public CC-BY~4.0 annotation file (credited);
code and derived artifacts carry the same terms. The misuse risk is the one
documented here -- an inadequately specified null model yields a confident but
wrong structural attribution -- hence our recommendation to report survival
counts.

\section{Conclusion}
\label{sec:conclusion}

We preregistered a nested null ladder to test whether connectome topology
confers a reproducible advantage on a structure-sensitive task. Controlled
low-in-degree graphs show a large advantage that a global-statistic ladder
cannot localise; a task-aware null recovers it. On real MaleCNS wiring a
fixed-density path-survival dose initially failed because
density-compensating swaps created fresh two-hop paths that global invariants
missed; a preregistered path-blocking repair yielded a monotone dose, and a
fully anticipatory replication on two independently constructed real subgraphs
reproduces it on both. This establishes the path-level mechanism on real
wiring for a controlled diagnostic label. It does not establish transfer to
natural visual tasks (\texttt{mechanism\_transfer\_confirmed=false};
\texttt{expected\_strong} unmet): a faithful rerun of the published ladder
reproduces its weak/strict pattern but not task-aware recovery, and a real
image task shows the wiring prior is not a competitive bias. Motif survival is
a measured predictor within this regime, not a universal cause. Our contribution is a
diagnostic standard for control sufficiency with machine-checked invariants:
control construction, not topology, explains conflicting reports.

\clearpage

\section*{LLM Usage Statement}
\emph{[To be supplied by the authors.]}

{
\small
\bibliographystyle{ieeenat_fullname}
\bibliography{main}
}

\end{document}
